\section{The second and third order terms}\label{cms:sec:explicit}
\subsection{The value and analytic slope at minus one}
At $z=-1$, consecutive product factors telescope:
\begin{equation}\label{cms:eq:minusvalue}
\left(1+\frac1{2k-1}\right)
\left(1-\frac1{2k}\right)=1.
\end{equation}
Thus $F(-1)=1$. The first terms of the filtered tree series are
\begin{equation}\label{cms:eq:B2}
B_2(s)=s-2s^2+O(s^3).
\end{equation}
To obtain the order-three coefficient, the analytic linear term of $\log F(-\ee^{-s})$ is also required. We compute it directly rather than infer it from finite coefficient data.

Isolate $j=1$, as before, and group the logarithmic sum by the power $r$. The analytic coefficient of $s$ in the $r=1$ term is
\[
\frac32-\log2,
\]
and in the $r=2$ term it is
\[
\zeta(2)-1-\log2.
\]
For clarity, the former comes from differentiating $z(\Li_1(z)-z)$ at $z=-\ee^{-s}$. For the latter, insert the expansion of $\Li_2(\ee^{-2s})$ into $-\tfrac12\ee^{-2s}(\Li_2(\ee^{-2s})-\ee^{-2s})$ and separate its $-s\log s$ part. The isolated factor $\log(1+\ee^{-2s})$ contributes $-1$.

For all $r\geq3$, the first derivative is absolutely summable after $j=1$ has been removed. Summing the geometric series in $r$ gives the remaining analytic slope
\begin{align}\label{cms:eq:slopeS}
S&=\sum_{j\geq2}(j+1)
\frac{\bigl((-1)^j/j\bigr)^3}{1-(-1)^j/j}\\
&=\sum_{\substack{j\geq2\\j\ \mathrm{even}}}
\left(\frac2{j-1}-\frac2j-\frac1{j^2}\right)
-\sum_{\substack{j\geq3\\j\ \mathrm{odd}}}\frac1{j^2}\notag\\
&=2\log2+1-\zeta(2).\notag
\end{align}
Adding the four contributions gives $1/2$. Therefore the local logarithm is
\begin{align}\label{cms:eq:logminus}
\log F(-\ee^{-s})={}&-s\log s+\frac{s}{2}+2s^2\log s\\
&+\text{analytic terms of degree at least }2
+O\bigl(s^3(1+|\log s|)\bigr),\notag
\end{align}
where the finite-jet statement and its differentiated control follow from the earlier factorization. The omitted analytic coefficient of $s^2$ will not affect the nonanalytic terms of degree at most two.

\subsection{Changing from the exponential coordinate}
Set $t=1+z$, so $s=-\log(1-t)=t+t^2/2+O(t^3)$. Exponentiating \eqref{cms:eq:logminus} first gives the nonanalytic terms
\[
-s\log s+s^2\left(\frac12\log^2s+\frac32\log s\right).
\]
The analytic $s/2$ term matters through its product with $-s\log s$. On changing coordinates, the contribution of $-s\log s$ includes $-\tfrac12t^2\log t$. The final nonanalytic jet through degree two is therefore
\begin{equation}\label{cms:eq:tminus}
-t\log t+t^2\left(\frac12\log^2t+\log t\right).
\end{equation}
This calculation explains why omitting the analytic linear coefficient, or forgetting the coordinate change, would alter the constant at order three.

\subsection{Exact transfer at minus one}
For $m\geq3$, the basic identities are
\begin{align}\label{cms:eq:exactlogidentities}
[z^m](1-z)\log(1-z)&=\frac1{m(m-1)},\\
[z^m](1-z)^2\log(1-z)&=-\frac2{m(m-1)(m-2)},\notag\\
[z^m](1-z)^2\log^2(1-z)&=
\frac{4(H_{m-3}-3/2)}{m(m-1)(m-2)}.\notag
\end{align}
The last identity follows either by twice differentiating \eqref{cms:eq:transfer} at $p=2$, or by multiplying the identity $[z^m]\log^2(1-z)=2H_{m-1}/m$ by $(1-z)^2$. Substituting $z\mapsto-z$ adds the factor $(-1)^m$.

Since $H_{m-3}=\log m+\gamma+O(m^{-1})$, the coefficient of \eqref{cms:eq:tminus} is
\begin{equation}\label{cms:eq:minusm}
(-1)^m\left\{-\frac1{m^2}
+\frac{2\log m+2\gamma-6}{m^3}\right\}
+O\!\left(\frac{(1+\log m)^2}{m^4}\right).
\end{equation}
For example, the three contributions to the constant at $m^{-3}$ are $-1$ from $-t\log t$, $-3$ from half the squared-log term, and $-2$ from $t^2\log t$. Their sum is $-6$.

Now $m=n-1$ and $(-1)^m=-(-1)^n$. Expanding the denominators and logarithm gives
\begin{equation}\label{cms:eq:minusn}
(-1)^n\left\{\frac1{n^2}
+\frac{8-2\gamma-2\log n}{n^3}\right\}
+O\!\left(\frac{(1+\log n)^2}{n^4}\right).
\end{equation}
The change from $6$ to $8$ is the $m=n-1$ shift of the leading term; this is not an optional convention.

\subsection{The two cubic roots}
For $\zeta=\om$ or $\overline\om$, equation \eqref{cms:eq:rootleading} becomes
\[
-\frac32F(\zeta)\left(1-\frac z\zeta\right)^2
\log\left(1-\frac z\zeta\right).
\]
By \eqref{cms:eq:singlelog}, its leading coefficient is
\[
3F(\zeta)\zeta^{-m}m^{-3}.
\]
The conjugate pair therefore contributes
\begin{equation}\label{cms:eq:cubicpair}
\frac{6\Rea(C_\om\om^{-m})}{m^3}
+O\!\left(\frac{1+\log m}{m^4}\right).
\end{equation}
Replacing $m$ by $n-1$ changes the phase to $\om^{1-n}$ and does not otherwise change this leading order. No roots of order at least four contribute through order three. Combining the pole coefficient $A$, \eqref{cms:eq:minusn}, and \eqref{cms:eq:cubicpair} proves Theorem~\ref{cms:thm:explicit}.

\subsection{The first periodic polynomials and numerical meaning}
The first two polynomials in Theorem~\ref{cms:thm:main} are explicitly
\begin{align*}
P_2(n,u)&=(-1)^n,\\
P_3(n,u)&=(-1)^n(8-2\gamma-2u)
 +6\Rea(C_\om\om^{1-n}).
\end{align*}
Their periods divide $2$ and $6$, respectively. The coefficient of $u$ and the constant coefficient in $P_3$ both average to zero over six consecutive integers. The absence of a smooth correction is visible at these first orders and is proved at every fixed order by the positive-root cancellation.

The code evaluates $C_\om$ from the Gamma product and independently from the polylogarithmic logarithm with $j=1$ removed. If the sum in that logarithm is truncated after $r=R$, its omitted absolute logarithmic tail is bounded by
\begin{equation}\label{cms:eq:numericaltail}
\sum_{r>R}\frac1r\sum_{j\geq2}j^{-r}
\leq\frac{3\,2^{-R}}{R+1}.
\end{equation}
This is an analytic truncation bound. Ordinary high precision evaluation of the retained transcendental terms still has numerical roundoff; the agreement of two evaluations is a diagnostic rather than a certified enclosure unless all those operations are evaluated with directed interval control. No coefficient in \eqref{cms:eq:explicit} is obtained by a numerical fit.
